\chapter{Απαιτήσεις και σχεδίαση του συστήματος}
\label{ch:apaitiseis}

Το κεφάλαιο περιγράφει το σύστημα πριν από την υλοποίησή του: ποιοι το χρησιμοποιούν, τι
οφείλει να κάνει, ποιες ιδιότητες πρέπει να διατηρεί, και ποιες σχεδιαστικές αποφάσεις
ελήφθησαν για να τα επιτύχει. Ακολουθείται η μεθοδολογία της \cite{hadjichristofi2022}:
ρόλοι, απαιτήσεις, περιπτώσεις χρήσης με διαγράμματα \en{UML} \cite{uml}, αρχιτεκτονική και
σχεδίαση διεπαφής.

\section{Η γενική ιδέα}
\label{sec:idea}

Το σύστημα είναι ένα \en{marketplace} στο οποίο δεν υπάρχουν τιμές. Κάθε χρήστης βλέπει τα
\en{NFTs} των υπολοίπων και, αν κάποιο τον ενδιαφέρει, υποβάλλει στον κάτοχό του μια
\emph{προσφορά}: ένα σύνολο δικών του \en{NFTs}, ένα ποσό \en{SUI}, ή συνδυασμό των δύο. Ο
κάτοχος βλέπει τις προσφορές που έχει λάβει και για καθεμία επιλέγει αν θα την αποδεχθεί ή
θα την απορρίψει. Η αποδοχή ολοκληρώνει την ανταλλαγή αμέσως. Όσο η προσφορά εκκρεμεί, τα
προσφερόμενα στοιχεία βρίσκονται δεσμευμένα στην αλυσίδα, και ο προσφέρων μπορεί να την
ακυρώσει οποτεδήποτε για να τα πάρει πίσω.

Η εφαρμογή επίδειξης περιλαμβάνει τρεις συλλογές (\en{Cars}, \en{Houses}, \en{Tickets}), η
καθεμία με τρεις βαθμίδες σπανιότητας (\en{common}, \en{uncommon}, \en{rare}), και επιτρέπει
σε οποιονδήποτε να δημιουργήσει δωρεάν \en{NFTs} για να πειραματισθεί. Το πρωτόκολλο όμως
δεν γνωρίζει τίποτε για τις συγκεκριμένες συλλογές: λειτουργεί με οποιονδήποτε τύπο
αντικειμένου.

\section{Ρόλοι χρηστών}
\label{sec:roles}

Το σύστημα δεν διακρίνει κατηγορίες λογαριασμών. Κάθε διεύθυνση μπορεί να εκτελέσει κάθε
λειτουργία, και οι ρόλοι που ακολουθούν ορίζονται ως προς μια συγκεκριμένη προσφορά.

\begin{description}[leftmargin=1.2em,itemsep=2pt]
  \item[Προσφέρων (\en{bidder}).] Ο χρήστης που υποβάλλει μια προσφορά. Είναι ο μόνος που
    μπορεί να την ακυρώσει.
  \item[Κάτοχος (\en{owner}).] Ο χρήστης που κατέχει το \en{NFT} στο οποίο στοχεύει η
    προσφορά. Μπορεί να την αποδεχθεί ή να την απορρίψει.
\end{description}

Ο ίδιος χρήστης είναι συνήθως προσφέρων σε κάποιες προσφορές και κάτοχος σε άλλες.
Αξίζει να τονισθεί η απουσία τρίτου ρόλου: δεν υπάρχει διαχειριστής. Το
Σχήμα~\ref{fig:usecase} δείχνει το διάγραμμα περιπτώσεων χρήσης.

\newcommand{\actor}[3]{%
  \draw (#1,#2+0.55) circle (0.18);
  \draw (#1,#2+0.37) -- (#1,#2-0.1) (#1-0.3,#2+0.2) -- (#1+0.3,#2+0.2)
        (#1,#2-0.1) -- (#1-0.25,#2-0.5) (#1,#2-0.1) -- (#1+0.25,#2-0.5);
  \node[font=\footnotesize, align=center] at (#1,#2-0.95) {#3};}

\begin{figure}[ht]
\centering
\begin{tikzpicture}[
  font=\footnotesize,
  uc/.style={draw, ellipse, minimum width=3.6cm, minimum height=0.8cm, align=center, fill=gray!8, inner sep=1pt}]
  \draw[rounded corners=4pt] (-4.6,-5.3) rectangle (4.6,3.6);
  \node[font=\small\bfseries] at (0,3.25) {\en{NFT Barter}};
  % common use cases
  \node[uc] (connect) at (0,2.4) {Σύνδεση πορτοφολιού};
  \node[uc] (mint) at (0,1.4) {Δημιουργία \en{NFT}};
  \node[uc] (browse) at (0,0.4) {Περιήγηση ανά συλλογή};
  % bidder
  \node[uc] (place) at (-2.2,-1.0) {Υποβολή προσφοράς};
  \node[uc] (placed) at (-2.2,-2.3) {Προβολή υποβληθεισών};
  \node[uc] (cancel) at (-2.2,-3.6) {Ακύρωση προσφοράς};
  % owner
  \node[uc] (recv) at (2.2,-1.7) {Προβολή εισερχόμενων};
  \node[uc] (accept) at (2.2,-3.0) {Αποδοχή προσφοράς};
  \node[uc] (reject) at (2.2,-4.3) {Απόρριψη προσφοράς};
  \actor{-6.2}{-1.6}{Προσφέρων};
  \actor{6.2}{-1.6}{Κάτοχος};
  \foreach \n in {connect,mint,browse,place,placed,cancel} \draw (-5.85,-1.4) -- (\n.west);
  \foreach \n in {connect,mint,browse,recv,accept,reject} \draw (5.85,-1.4) -- (\n.east);
\end{tikzpicture}
\caption{Διάγραμμα περιπτώσεων χρήσης. Οι δύο ρόλοι ορίζονται ανά προσφορά· ο ίδιος χρήστης
αναλαμβάνει και τους δύο.}
\label{fig:usecase}
\end{figure}

\section{Λειτουργικές απαιτήσεις}
\label{sec:functional}

Οι λειτουργικές απαιτήσεις συγκεντρώνονται στον Πίνακα~\ref{tab:functional}. Οι ΛΑ1 έως
ΛΑ6 αφορούν το πρωτόκολλο και υλοποιούνται στο έξυπνο συμβόλαιο· οι ΛΑ7 έως ΛΑ10 αφορούν
την εφαρμογή.

\begin{table}[H]
\centering
\caption{Λειτουργικές απαιτήσεις}
\label{tab:functional}
\small
\begin{tabular}{lp{12.6cm}}
\toprule
Κωδ. & Απαίτηση \\
\midrule
ΛΑ1 & Ο χρήστης δημιουργεί (\en{mint}) \en{NFT} με όνομα, περιγραφή, εικόνα, συλλογή και σπανιότητα. \\
ΛΑ2 & Ο χρήστης υποβάλλει προσφορά που περιέχει μηδέν ή περισσότερα \en{NFTs} ενός τύπου και μηδέν ή περισσότερο \en{SUI}, με τουλάχιστον ένα από τα δύο μη μηδενικό. \\
ΛΑ3 & Η προσφορά μπορεί να στοχεύει συγκεκριμένο \en{NFT}, οπότε μόνο με αυτό μπορεί να γίνει αποδεκτή. \\
ΛΑ4 & Ο κάτοχος αποδέχεται προσφορά παραδίδοντας τα ζητούμενα \en{NFTs}· λαμβάνει όλα τα προσφερόμενα στοιχεία και ο προσφέρων λαμβάνει τα ζητούμενα, στην ίδια συναλλαγή. \\
ΛΑ5 & Ο κάτοχος του στοχευόμενου \en{NFT} απορρίπτει προσφορά· όλα τα προσφερόμενα στοιχεία επιστρέφουν στον προσφέροντα. \\
ΛΑ6 & Ο προσφέρων ακυρώνει προσφορά του· όλα τα προσφερόμενα στοιχεία επιστρέφουν σε αυτόν. \\
\midrule
ΛΑ7 & Η εφαρμογή παρουσιάζει τα διαθέσιμα \en{NFTs} ανά συλλογή, με κάτοχο και σπανιότητα, και επιτρέπει φιλτράρισμα ανά βαθμίδα σπανιότητας. \\
ΛΑ8 & Η εφαρμογή παρουσιάζει στον χρήστη τα \en{NFTs} που κατέχει και, για καθένα, τις προσφορές που έχει λάβει, καθώς και όσα έχει δεσμεύσει σε δικές του προσφορές. \\
ΛΑ9 & Η εφαρμογή παρουσιάζει στον χρήστη τις προσφορές που έχει υποβάλει και εκκρεμούν. \\
ΛΑ10 & Για κάθε προσφορά η εφαρμογή δείχνει ρητά τι δίνει και τι λαμβάνει ο χρήστης. \\
\bottomrule
\end{tabular}
\end{table}

\section{Μη λειτουργικές απαιτήσεις}
\label{sec:nonfunctional}

Όπως παρατηρεί και η \cite{hadjichristofi2022}, σε εφαρμογές αυτού του είδους είναι οι μη
λειτουργικές απαιτήσεις που καθορίζουν την αρχιτεκτονική. Συγκεντρώνονται στον
Πίνακα~\ref{tab:nonfunctional}.

\begin{table}[H]
\centering
\caption{Μη λειτουργικές απαιτήσεις}
\label{tab:nonfunctional}
\small
\begin{tabular}{lp{3.0cm}p{9.4cm}}
\toprule
Κωδ. & Ιδιότητα & Απαίτηση \\
\midrule
ΜΛΑ1 & Ατομικότητα & Η ανταλλαγή είτε ολοκληρώνεται πλήρως και για τα δύο μέρη είτε δεν συμβαίνει καθόλου. \\
ΜΛΑ2 & Διατήρηση στοιχείων & Κάθε δεσμευμένο στοιχείο καταλήγει είτε στον αντισυμβαλλόμενο (αποδοχή) είτε στον προσφέροντα (απόρριψη, ακύρωση). Δεν υπάρχει διαδρομή στην οποία χάνεται ή παραμένει εγκλωβισμένο. \\
ΜΛΑ3 & Εξουσιοδότηση & Μόνο ο προσφέρων ακυρώνει· μόνο ο κάτοχος του στοχευόμενου \en{NFT} αποδέχεται ή απορρίπτει. \\
ΜΛΑ4 & Απουσία έμπιστου τρίτου & Δεν υπάρχει θεματοφύλακας, διαχειριστής ή προνομιακό κλειδί, ούτε τέλη πέραν του \en{gas}. \\
ΜΛΑ5 & Γενικότητα & Το πρωτόκολλο λειτουργεί με οποιονδήποτε τύπο αντικειμένου, χωρίς καταχώρηση συλλογών. \\
ΜΛΑ6 & Ορθότητα υπό ταυτόχρονη πρόσβαση & Ταυτόχρονες ενέργειες στην ίδια προσφορά (π.χ. αποδοχή και ακύρωση) έχουν καλά ορισμένο αποτέλεσμα: εκτελείται ακριβώς μία. \\
ΜΛΑ7 & Διαφάνεια & Η κατάσταση κάθε προσφοράς και τα δεσμευμένα στοιχεία είναι επαληθεύσιμα από οποιονδήποτε με δεδομένα της αλυσίδας. \\
ΜΛΑ8 & Ευχρηστία & Κάθε ενέργεια απαιτεί μία υπογραφή. Η εφαρμογή λειτουργεί με κάθε πορτοφόλι του \en{Wallet Standard}, σε επιτραπέζια και κινητή οθόνη. \\
\bottomrule
\end{tabular}
\end{table}

\section{Κύκλος ζωής της προσφοράς}
\label{sec:lifecycle}

Μια προσφορά έχει μία μόνο ενεργή κατάσταση και τρεις τερματικές, όπως δείχνει το
Σχήμα~\ref{fig:state}. Δεν υπάρχουν ενδιάμεσες καταστάσεις, διότι κάθε μετάβαση είναι μία
ατομική συναλλαγή.

\begin{figure}[ht]
\centering
\begin{tikzpicture}[
  font=\small, >=Stealth,
  st/.style={draw, rounded corners=8pt, minimum width=2.9cm, minimum height=1.0cm, align=center, fill=gray!8},
  ev/.style={font=\footnotesize, align=center}]
  \node[circle, fill=black, minimum size=0.28cm, inner sep=0] (init) at (-4.2,0) {};
  \node[st] (placed) at (0,0) {\textbf{Ενεργή}\\{\footnotesize στοιχεία δεσμευμένα}};
  \node[st] (acc) at (6.0,2.2) {\textbf{Αποδεκτή}\\{\footnotesize ανταλλαγή ολοκληρώθηκε}};
  \node[st] (rej) at (6.0,0) {\textbf{Απορριφθείσα}\\{\footnotesize επιστροφή στον προσφέροντα}};
  \node[st] (can) at (6.0,-2.2) {\textbf{Ακυρωμένη}\\{\footnotesize επιστροφή στον προσφέροντα}};
  \draw[->] (init) -- node[ev, above] {\code{place\_bid}} node[ev, below] {προσφέρων} (placed);
  \draw[->] (placed.north east) -- node[ev, above, sloped] {\code{accept\_bid}} node[ev, below, sloped] {κάτοχος} (acc.west);
  \draw[->] (placed) -- node[ev, above] {\code{reject\_bid}} node[ev, below] {κάτοχος} (rej);
  \draw[->] (placed.south east) -- node[ev, above, sloped] {\code{cancel\_bid}} node[ev, below, sloped] {προσφέρων} (can.west);
\end{tikzpicture}
\caption{Μηχανή καταστάσεων της προσφοράς. Σε κάθε τερματική κατάσταση το αντικείμενο
\en{Bid} έχει διαγραφεί από την αλυσίδα.}
\label{fig:state}
\end{figure}

Ο κύκλος ζωής συνοδεύεται από τρεις αναλλοίωτες συνθήκες, τις οποίες η υλοποίηση οφείλει
να διατηρεί και οι οποίες θα ελεγχθούν στο Κεφάλαιο~\ref{ch:contract}:

\begin{description}[leftmargin=1.2em,itemsep=2pt]
  \item[Α1.] Όσο η προσφορά είναι ενεργή, τα δεσμευμένα στοιχεία ανήκουν στο αντικείμενο
    της προσφοράς και σε κανέναν χρήστη.
  \item[Α2.] Κάθε μετάβαση σε τερματική κατάσταση αδειάζει πλήρως τη δέσμευση και διαγράφει
    το αντικείμενο της προσφοράς. Καμία προσφορά δεν μπορεί να εκτελεστεί δύο φορές.
  \item[Α3.] Τα στοιχεία που εξέρχονται από τη δέσμευση είναι ακριβώς εκείνα που εισήλθαν.
\end{description}

\section{Σχεδιαστικές αποφάσεις}
\label{sec:decisions}

Οι απαιτήσεις επιδέχονται περισσότερες από μία υλοποιήσεις. Η ενότητα καταγράφει τις
κύριες αποφάσεις, τις εναλλακτικές που εξετάσθηκαν και τους λόγους της επιλογής.

\subsection{Η προσφορά ως κοινόχρηστο αντικείμενο}

Η προσφορά πρέπει να είναι προσβάσιμη από δύο διαφορετικούς χρήστες: τον προσφέροντα, για
την ακύρωση, και τον κάτοχο, για την αποδοχή ή την απόρριψη. Αν ήταν αντικείμενο του
προσφέροντος, ο κάτοχος δεν θα μπορούσε να τη χρησιμοποιήσει σε συναλλαγή του. Αν
μεταβιβαζόταν στον κάτοχο, ο προσφέρων θα έχανε τη δυνατότητα ακύρωσης και θα εξαρτιόταν
από την καλή θέληση του αποδέκτη για να πάρει πίσω τα στοιχεία του, κάτι που παραβιάζει τη
ΜΛΑ4. Η μόνη μορφή ιδιοκτησίας που ικανοποιεί και τους δύο είναι το κοινόχρηστο
αντικείμενο, με τους ελέγχους εξουσιοδότησης να γίνονται από τον κώδικα.

Εξετάσθηκε και η εναλλακτική ενός ενιαίου κοινόχρηστου \emph{μητρώου} που θα περιείχε όλες
τις προσφορές. Απορρίφθηκε διότι κάθε συναλλαγή του συστήματος θα έπρεπε τότε να κλειδώσει
το ίδιο αντικείμενο, και άσχετες μεταξύ τους προσφορές θα ανταγωνίζονταν για διάταξη. Με
ένα κοινόχρηστο αντικείμενο \emph{ανά προσφορά}, ανταγωνισμός υπάρχει μόνο μεταξύ
ενεργειών στην ίδια προσφορά, δηλαδή ακριβώς εκεί όπου η διάταξη είναι απαραίτητη (ΜΛΑ6).
Το κόστος της επιλογής είναι ότι οι προσφορές δεν είναι απαριθμήσιμες από κάποιο κεντρικό
σημείο· το ζήτημα αντιμετωπίζεται με συμβάντα, όπως εξηγείται παρακάτω.

\subsection{Δέσμευση με δυναμικά πεδία αντικειμένων}

Τα προσφερόμενα \en{NFTs} πρέπει να αποθηκευθούν μέσα στην προσφορά. Η απλούστερη λύση,
ένα πεδίο \code{vector<Offered>}, θα περιτύλιγε τα αντικείμενα: θα έπαυαν να εμφανίζονται
ως αυτοτελή αντικείμενα και ένας τρίτος δεν θα μπορούσε εύκολα να επαληθεύσει τι περιέχει
η προσφορά. Προτιμήθηκαν τα δυναμικά πεδία αντικειμένων
(Ενότητα~\ref{sec:dynamic-fields}), με τα οποία κάθε δεσμευμένο \en{NFT} διατηρεί το
\en{ID} του και εμφανίζεται στους εξερευνητές της αλυσίδας ως παιδί της προσφοράς (ΜΛΑ7). Ως κλειδί
χρησιμοποιείται ένας αύξων αριθμός, περιτυλιγμένος σε ειδικό τύπο ώστε να μην συγκρούεται
με πεδία που θα μπορούσε να προσθέσει άλλος κώδικας.

\subsection{Γενικοί τύποι για τις συλλογές}

Η υποστήριξη διαφορετικών συλλογών (ΜΛΑ5) θα μπορούσε να επιτευχθεί με ένα μητρώο
επιτρεπόμενων τύπων, το οποίο όμως προϋποθέτει διαχειριστή. Επιλέχθηκε η παραμετροποίηση
της προσφοράς με δύο παραμέτρους τύπου, \code{Bid<Offered, Requested>}. Ο προσφέρων
καθορίζει και τους δύο κατά την υποβολή, και οι συναρτήσεις αποδοχής, απόρριψης και
ακύρωσης απαιτούν τους ίδιους τύπους για να δεχθούν το αντικείμενο. Ο έλεγχος γίνεται από
την εικονική μηχανή: είναι αδύνατο να γίνει αποδεκτή μια προσφορά για
\code{Requested = Shield} παραδίδοντας αντικείμενο άλλου τύπου, διότι η κλήση δεν θα
περάσει καν τον έλεγχο τύπων.

Η επιλογή έχει έναν περιορισμό που πρέπει να δηλωθεί: τα προσφερόμενα \en{NFTs} μιας
προσφοράς έχουν όλα τον ίδιο τύπο \code{Offered}. Ο προσφέρων μπορεί να δώσει αντικείμενα
τύπου Α για ένα αντικείμενο τύπου Β, όχι όμως αντικείμενα τύπων Α και Γ μαζί στην ίδια
προσφορά. Η άρση του περιορισμού συζητείται στο Κεφάλαιο~\ref{ch:symperasmata}.

\subsection{Στοχευμένες προσφορές}

Μια προσφορά μπορεί να φέρει προαιρετικά το \en{ID} του \en{NFT} που ζητά. Όταν το φέρει,
η αποδοχή απαιτεί να παραδοθεί ακριβώς αυτό το αντικείμενο. Όταν δεν το φέρει, η προσφορά
είναι \emph{ανοικτή}: μπορεί να την αποδεχθεί οποιοσδήποτε παραδώσει το ζητούμενο πλήθος
αντικειμένων του τύπου \code{Requested}. Το συμβόλαιο υποστηρίζει και τις δύο μορφές. Η
εφαρμογή χρησιμοποιεί αποκλειστικά τη στοχευμένη, που αντιστοιχεί στο σενάριο «θέλω αυτό
το συγκεκριμένο αντικείμενο».

Η απόρριψη έχει νόημα μόνο για στοχευμένες προσφορές, αφού μια ανοικτή δεν απευθύνεται σε
κανέναν συγκεκριμένα. Η εξουσιοδότησή της θέτει ένα ενδιαφέρον ερώτημα: πώς αποδεικνύει ο
καλών ότι κατέχει το στοχευόμενο \en{NFT}; Στο \en{Sui} δεν χρειάζεται να ερωτηθεί κάποιο
μητρώο. Αρκεί η συνάρτηση να ζητήσει ως όρισμα μια \emph{αναφορά} στο αντικείμενο: μόνο ο
κάτοχος ενός αντικειμένου διεύθυνσης μπορεί να το περάσει ως είσοδο σε συναλλαγή, άρα η
ίδια η ικανότητα να κληθεί η συνάρτηση αποτελεί την απόδειξη.

Αξίζει επίσης να σημειωθεί ότι η προσφορά στοχεύει το \emph{αντικείμενο} και όχι τον
κάτοχό του. Αν το \en{NFT} αλλάξει χέρια όσο η προσφορά εκκρεμεί, ο νέος κάτοχος είναι
εκείνος που μπορεί πλέον να την αποδεχθεί ή να την απορρίψει.

\subsection{Το νόμισμα ως μέρος της προσφοράς}

Το προσφερόμενο \en{SUI} αποθηκεύεται ως \code{Balance<SUI>} σε πεδίο της προσφοράς. Η
συνάρτηση υποβολής δέχεται πάντοτε ένα \code{Coin<SUI>}, ενδεχομένως μηδενικής αξίας, ώστε
να υπάρχει μία υπογραφή συνάρτησης για όλους τους συνδυασμούς. Η μόνη απαίτηση είναι να
προσφέρεται \emph{κάτι}: τουλάχιστον ένα \en{NFT} ή θετικό ποσό.

\subsection{Ανακάλυψη μέσω συμβάντων}

Εφόσον δεν υπάρχει μητρώο προσφορών και ένας πλήρης κόμβος δεν απαριθμεί κοινόχρηστα
αντικείμενα ανά τύπο (Ενότητα~\ref{sec:events}), η εφαρμογή χρειάζεται άλλον τρόπο να βρει
τις προσφορές. Το συμβόλαιο εκπέμπει ένα συμβάν για κάθε μετάβαση του κύκλου ζωής, καθώς
και για κάθε δημιουργία \en{NFT}. Η εφαρμογή διαβάζει τα συμβάντα υποβολής, συλλέγει τα
\en{IDs} των προσφορών, και ανακτά τα αντίστοιχα αντικείμενα· όσα δεν υπάρχουν πλέον έχουν
φτάσει σε τερματική κατάσταση. Η προσέγγιση δεν απαιτεί καμία υποδομή πέραν ενός
πλήρους κόμβου. Τα όριά της
ως προς την κλιμάκωση συζητούνται στην Ενότητα~\ref{sec:limitations}.

\section{Περιπτώσεις χρήσης}
\label{sec:usecases}

\newcommand{\ucrow}[2]{\textbf{#1} & #2 \\[2pt]}
\newenvironment{usecase}[1]{%
  \par\medskip\noindent\begin{tabular}{@{}p{3.1cm}p{11.6cm}@{}}
  \toprule
  \multicolumn{2}{@{}l}{\textbf{#1}} \\
  \midrule}{%
  \bottomrule
  \end{tabular}\par\medskip}

\subsection{ΠΧ1: Σύνδεση πορτοφολιού}

\begin{usecase}{ΠΧ1 — Σύνδεση πορτοφολιού}
\ucrow{Ηθοποιός}{Οποιοσδήποτε χρήστης}
\ucrow{Προϋποθέσεις}{Υπάρχει εγκατεστημένο πορτοφόλι συμβατό με το \en{Wallet Standard}.}
\ucrow{Κύρια ροή}{1.~Ο χρήστης επιλέγει «\en{Connect Wallet}». 2.~Η εφαρμογή εμφανίζει τα διαθέσιμα πορτοφόλια. 3.~Ο χρήστης επιλέγει ένα και εγκρίνει τη σύνδεση. 4.~Η εφαρμογή εμφανίζει τη διεύθυνσή του και ενεργοποιεί τις λειτουργίες που απαιτούν υπογραφή.}
\ucrow{Εναλλακτικές}{Ο χρήστης απορρίπτει τη σύνδεση: η εφαρμογή παραμένει σε κατάσταση ανάγνωσης.}
\ucrow{Μετασυνθήκες}{Η επιλογή αποθηκεύεται τοπικά, ώστε η σύνδεση να αποκατασταθεί αυτόματα στην επόμενη επίσκεψη.}
\end{usecase}

Χωρίς συνδεδεμένο πορτοφόλι ο χρήστης μπορεί να περιηγηθεί στις συλλογές, όχι όμως να
υποβάλει προσφορά ή να δει τις σελίδες που αφορούν δικά του στοιχεία.

\subsection{ΠΧ2: Δημιουργία NFT}

\begin{usecase}{ΠΧ2 — Δημιουργία \en{NFT}}
\ucrow{Ηθοποιός}{Συνδεδεμένος χρήστης}
\ucrow{Κύρια ροή}{1.~Ο χρήστης επιλέγει ένα από τα εννέα είδη του καταλόγου και ποσότητα. 2.~Η εφαρμογή συνθέτει ένα \en{PTB} με μία κλήση \code{mint} ανά τεμάχιο και μία τελική εντολή που μεταβιβάζει τα αποτελέσματα στη διεύθυνσή του. 3.~Ο χρήστης υπογράφει. 4.~Τα νέα \en{NFTs} βρίσκονται στο πορτοφόλι του.}
\ucrow{Εναλλακτικές}{Ο χρήστης απορρίπτει την υπογραφή ή δεν έχει \en{SUI} για \en{gas}: εμφανίζεται μήνυμα αποτυχίας και δεν αλλάζει τίποτε. Σε δίκτυο δοκιμών μπορεί να ζητήσει \en{SUI} από το κουμπί «\en{Faucet}» της εφαρμογής. Αν ζητήσει δεύτερη δημιουργία ενώ η πρώτη εκκρεμεί, η δεύτερη μπαίνει σε ουρά και εκτελείται μετά την πρώτη.}
\ucrow{Μετασυνθήκες}{Εκπέμπεται συμβάν \code{NftMinted} ανά τεμάχιο. Τα \en{NFTs} εμφανίζονται στη σελίδα της συλλογής τους.}
\end{usecase}

\subsection{ΠΧ3: Υποβολή προσφοράς}

\begin{usecase}{ΠΧ3 — Υποβολή προσφοράς}
\ucrow{Ηθοποιός}{Προσφέρων}
\ucrow{Προϋποθέσεις}{Συνδεδεμένο πορτοφόλι. Το \en{NFT}-στόχος ανήκει σε άλλον χρήστη.}
\ucrow{Κύρια ροή}{1.~Ο χρήστης επιλέγει «\en{Place a bid}» σε ένα \en{NFT}. 2.~Ορίζει ποσό \en{SUI} ή παραλείπει το βήμα («\en{Skip}»). 3.~Επιλέγει προαιρετικά \en{NFTs} από το πορτοφόλι του. 4.~Ελέγχει τη σύνοψη «δίνεις / παίρνεις» και επιβεβαιώνει. 5.~Η εφαρμογή συνθέτει το \en{PTB} του Σχήματος~\ref{fig:ptb} και ο χρήστης υπογράφει. 6.~Δημιουργείται κοινόχρηστο αντικείμενο \code{Bid} με τα στοιχεία δεσμευμένα.}
\ucrow{Εναλλακτικές}{3α.~Δεν έχει επιλεγεί ούτε \en{SUI} ούτε \en{NFT}: η εφαρμογή δεν επιτρέπει τη μετάβαση στην επισκόπηση. 5α.~Ο χρήστης απορρίπτει την υπογραφή: δεν δεσμεύεται τίποτε.}
\ucrow{Μετασυνθήκες}{Τα προσφερόμενα στοιχεία δεν βρίσκονται πλέον στο πορτοφόλι του προσφέροντος. Εκπέμπεται \code{BidPlaced}. Η προσφορά εμφανίζεται στον κάτοχο ως εισερχόμενη.}
\end{usecase}

Το Σχήμα~\ref{fig:seq-place} δείχνει τη ροή μηνυμάτων. Αξίζει να παρατηρηθεί ότι η
εφαρμογή δεν αποστέλλει η ίδια τη συναλλαγή: την παραδίδει στο πορτοφόλι, το οποίο την
παρουσιάζει στον χρήστη, την υπογράφει και την υποβάλλει.

\begin{figure}[ht]
\centering
\begin{tikzpicture}[font=\footnotesize, >=Stealth, lab/.style={above, fill=white, inner sep=1.5pt, outer sep=1pt}]
  \def\xa{0}\def\xb{3.3}\def\xc{6.6}\def\xd{9.9}\def\xe{13.0}
  \foreach \x/\n in {\xa/Προσφέρων,\xb/\en{dApp},\xc/Πορτοφόλι,\xd/Πλήρης κόμβος,\xe/\code{barter}} {
    \node[draw, fill=gray!8, minimum width=2.2cm, minimum height=0.6cm] at (\x,0) {\n};
    \draw[densely dashed] (\x,-0.3) -- (\x,-8.6);
  }
  \draw[->] (\xa,-0.9) -- node[lab] {ποσό \en{SUI}, \en{NFTs}} (\xb,-0.9);
  \draw[->] (\xb,-1.6) -- node[lab] {ανάγνωση \en{NFTs} χρήστη} (\xd,-1.6);
  \draw[<-, densely dashed] (\xb,-2.2) -- node[lab] {αντικείμενα} (\xd,-2.2);
  \draw[->] (\xa,-2.9) -- node[lab] {επιβεβαίωση} (\xb,-2.9);
  \draw[->] (\xb,-3.6) -- node[lab] {\en{PTB} προς υπογραφή} (\xc,-3.6);
  \draw[->] (\xc,-4.3) -- node[lab] {αίτημα έγκρισης} (\xa,-4.3);
  \draw[<-] (\xc,-4.9) -- node[lab] {έγκριση} (\xa,-4.9);
  \draw[->] (\xc,-5.6) -- node[lab] {υπογεγραμμένη συναλλαγή} (\xd,-5.6);
  \draw[->] (\xd,-6.3) -- node[lab] {\code{place\_bid}} (\xe,-6.3);
  \draw[<-, densely dashed] (\xd,-6.9) -- node[lab] {\code{Bid}, \code{BidPlaced}} (\xe,-6.9);
  \draw[<-, densely dashed] (\xb,-7.6) -- node[lab] {αποτέλεσμα συναλλαγής} (\xd,-7.6);
  \draw[<-, densely dashed] (\xa,-8.3) -- node[lab] {«\en{Bid placed}»} (\xb,-8.3);
\end{tikzpicture}
\caption{Διάγραμμα ακολουθίας της υποβολής προσφοράς.}
\label{fig:seq-place}
\end{figure}

\subsection{ΠΧ4: Αποδοχή προσφοράς}

\begin{usecase}{ΠΧ4 — Αποδοχή προσφοράς}
\ucrow{Ηθοποιός}{Κάτοχος}
\ucrow{Προϋποθέσεις}{Υπάρχει ενεργή προσφορά που στοχεύει \en{NFT} του χρήστη.}
\ucrow{Κύρια ροή}{1.~Ο χρήστης ανοίγει τις εισερχόμενες προσφορές, από το \en{NFT} στη σελίδα «\en{My NFTs}» ή από την καρτέλα «\en{Received}» της σελίδας «\en{My bids}». 2.~Εξετάζει τι δίνει και τι λαμβάνει. 3.~Επιλέγει «\en{Accept trade}» και υπογράφει. 4.~Στην ίδια συναλλαγή το \en{NFT} του μεταβιβάζεται στον προσφέροντα και τα δεσμευμένα στοιχεία σε αυτόν.}
\ucrow{Εναλλακτικές}{3α.~Η προσφορά ακυρώθηκε στο μεταξύ από τον προσφέροντα: η συναλλαγή απορρίπτεται πριν την υπογραφή, εμφανίζεται μήνυμα αποτυχίας, και η λίστα ανανεώνεται στην επόμενη ανάγνωση. 3β.~Το στοχευόμενο \en{NFT} άλλαξε κάτοχο: η προσφορά παύει να εμφανίζεται στον παλιό κάτοχο και εμφανίζεται στον νέο.}
\ucrow{Μετασυνθήκες}{Το αντικείμενο \code{Bid} έχει διαγραφεί. Εκπέμπεται \code{BidAccepted}.}
\end{usecase}

\begin{figure}[ht]
\centering
\begin{tikzpicture}[font=\footnotesize, >=Stealth, lab/.style={above, fill=white, inner sep=1.5pt, outer sep=1pt}]
  \def\xa{0}\def\xb{3.3}\def\xc{6.6}\def\xd{9.9}\def\xe{13.0}
  \foreach \x/\n in {\xa/Κάτοχος,\xb/\en{dApp},\xc/Πορτοφόλι,\xd/Πλήρης κόμβος,\xe/\code{barter}} {
    \node[draw, fill=gray!8, minimum width=2.2cm, minimum height=0.6cm] at (\x,0) {\n};
    \draw[densely dashed] (\x,-0.3) -- (\x,-8.6);
  }
  \draw[->] (\xb,-0.9) -- node[lab] {συμβάντα \code{BidPlaced}} (\xd,-0.9);
  \draw[<-, densely dashed] (\xb,-1.5) -- node[lab] {\en{IDs} προσφορών} (\xd,-1.5);
  \draw[->] (\xb,-2.2) -- node[lab] {αντικείμενα \code{Bid}, δεσμευμένα \en{NFTs}} (\xd,-2.2);
  \draw[<-, densely dashed] (\xa,-2.9) -- node[lab] {«δίνεις / παίρνεις»} (\xb,-2.9);
  \draw[->] (\xa,-3.6) -- node[lab] {\en{Accept trade}} (\xb,-3.6);
  \draw[->] (\xb,-4.3) -- node[lab] {\en{PTB} προς υπογραφή} (\xc,-4.3);
  \draw[<->] (\xc,-5.0) -- node[lab] {έγκριση} (\xa,-5.0);
  \draw[->] (\xc,-5.7) -- node[lab] {υπογεγραμμένη συναλλαγή} (\xd,-5.7);
  \draw[->] (\xd,-6.4) -- node[lab] {\code{accept\_bid}} (\xe,-6.4);
  \draw[<-, densely dashed] (\xd,-7.0) -- node[lab] {\code{BidAccepted}} (\xe,-7.0);
  \draw[<-, densely dashed] (\xb,-7.7) -- node[lab] {αποτέλεσμα συναλλαγής} (\xd,-7.7);
  \draw[<-, densely dashed] (\xa,-8.3) -- node[lab] {«\en{Trade complete}»} (\xb,-8.3);
\end{tikzpicture}
\caption{Διάγραμμα ακολουθίας της αποδοχής. Τα τρία πρώτα μηνύματα δείχνουν πώς η εφαρμογή
ανακαλύπτει τις εισερχόμενες προσφορές χωρίς \en{indexer}.}
\label{fig:seq-accept}
\end{figure}

\subsection{ΠΧ5: Απόρριψη προσφοράς}

\begin{usecase}{ΠΧ5 — Απόρριψη προσφοράς}
\ucrow{Ηθοποιός}{Κάτοχος}
\ucrow{Προϋποθέσεις}{Υπάρχει ενεργή στοχευμένη προσφορά για \en{NFT} του χρήστη.}
\ucrow{Κύρια ροή}{1.~Ο χρήστης επιλέγει «\en{Reject}» σε μια εισερχόμενη προσφορά και υπογράφει. 2.~Το συμβόλαιο επαληθεύει ότι ο καλών παρουσιάζει το στοχευόμενο \en{NFT}. 3.~Όλα τα δεσμευμένα στοιχεία επιστρέφουν στον προσφέροντα.}
\ucrow{Μετασυνθήκες}{Το \en{NFT} του κατόχου δεν μετακινείται. Το \code{Bid} διαγράφεται. Εκπέμπεται \code{BidRejected}.}
\end{usecase}

Η απόρριψη δεν είναι απαραίτητη για την ορθότητα· ο κάτοχος θα μπορούσε απλώς να αγνοήσει
μια προσφορά. Προσφέρει όμως δύο πράγματα: καθαρίζει τη λίστα του κατόχου, και
αποδεσμεύει άμεσα τα στοιχεία του προσφέροντος χωρίς να χρειαστεί εκείνος να ενεργήσει.

\subsection{ΠΧ6: Ακύρωση προσφοράς}

\begin{usecase}{ΠΧ6 — Ακύρωση προσφοράς}
\ucrow{Ηθοποιός}{Προσφέρων}
\ucrow{Προϋποθέσεις}{Ο χρήστης έχει υποβάλει προσφορά που είναι ακόμη ενεργή.}
\ucrow{Κύρια ροή}{1.~Ο χρήστης ανοίγει την καρτέλα «\en{Placed}» της σελίδας «\en{My bids}». 2.~Επιλέγει «\en{Cancel bid}» και υπογράφει. 3.~Όλα τα δεσμευμένα στοιχεία επιστρέφουν στο πορτοφόλι του.}
\ucrow{Εναλλακτικές}{2α.~Η προσφορά έγινε αποδεκτή στο μεταξύ: η συναλλαγή απορρίπτεται πριν την υπογραφή, εμφανίζεται μήνυμα αποτυχίας, και η λίστα ανανεώνεται στην επόμενη ανάγνωση.}
\ucrow{Μετασυνθήκες}{Το \code{Bid} διαγράφεται. Εκπέμπεται \code{BidCancelled}.}
\end{usecase}

\section{Αρχιτεκτονική}
\label{sec:architecture}

Το σύστημα αποτελείται από δύο μόνο μέρη που αναπτύχθηκαν στο πλαίσιο της εργασίας: το
πακέτο \en{Move} που δημοσιεύεται στην αλυσίδα, και την εφαρμογή ιστού που εκτελείται
εξ ολοκλήρου στον φυλλομετρητή. Δεν υπάρχει διακομιστής εφαρμογής ούτε βάση δεδομένων. Το
Σχήμα~\ref{fig:architecture} δείχνει τα μέρη και τις μεταξύ τους εξαρτήσεις.

\begin{figure}[ht]
\centering
\begin{tikzpicture}[
  font=\footnotesize, >=Stealth,
  box/.style={draw, rounded corners=3pt, minimum height=0.75cm, align=center, fill=gray!8},
  grp/.style={draw, rounded corners=5pt, dashed, inner sep=8pt}]
  % browser
  \node[box, minimum width=5.2cm] (ui) at (0,0) {Διεπαφή (\en{React}, \en{Tailwind})};
  \node[box, minimum width=5.2cm] (hooks) at (0,-1.1) {\en{Hooks} + \en{TanStack Query}};
  \node[box, minimum width=2.5cm] (kit) at (-1.35,-2.2) {\en{dApp Kit}};
  \node[box, minimum width=2.5cm] (sdk) at (1.35,-2.2) {\en{Sui TS SDK}};
  \node[grp, fit=(ui)(hooks)(kit)(sdk), label={[font=\footnotesize\bfseries]above:Εφαρμογή ιστού (φυλλομετρητής)}] (app) {};
  \node[box, minimum width=2.6cm] (wallet) at (-1.35,-4.3) {Πορτοφόλι};
  % network
  \node[box, minimum width=3.4cm] (node) at (8.6,-0.3) {Πλήρης κόμβος\\{\scriptsize\en{JSON-RPC} / \en{gRPC}}};
  \node[box, minimum width=3.4cm] (val) at (8.6,-2.0) {Επικυρωτές};
  \node[box, minimum width=3.4cm] (pkg) at (8.6,-3.7) {Πακέτο \code{nft\_barter}\\{\scriptsize\code{barter}, \code{marketplace\_nft}}};
  \node[grp, fit=(node)(val)(pkg), label={[font=\footnotesize\bfseries]above:Δίκτυο \en{Sui}}] (net) {};
  \draw[->] (ui) -- (hooks);
  \draw[->] (hooks) -- (kit);
  \draw[->] (hooks) -- (sdk);
  \draw[<->] (kit) -- node[left] {\en{Wallet Standard}} (wallet);
  \draw[->] (sdk.east) -- ++(1.0,0) |- node[pos=0.78, above] {αναγνώσεις} ([yshift=0.12cm]node.west);
  \draw[->] (wallet.east) -- node[below, pos=0.5] {υπογεγραμμένες συναλλαγές} ++(5.2,0) |- ([yshift=-0.15cm]node.west);
  \draw[<->] (node) -- (val);
  \draw[->] (val) -- node[right] {εκτέλεση} (pkg);
\end{tikzpicture}
\caption{Αρχιτεκτονική του συστήματος.}
\label{fig:architecture}
\end{figure}

Η εφαρμογή επικοινωνεί με το δίκτυο με δύο τρόπους. Οι \emph{αναγνώσεις} (αντικείμενα,
συμβάντα) γίνονται απευθείας από τον φυλλομετρητή προς έναν πλήρη κόμβο, μέσω του
\en{Sui TypeScript SDK}. Οι \emph{εγγραφές} περνούν πάντοτε από το πορτοφόλι: η εφαρμογή
κατασκευάζει τη συναλλαγή, το πορτοφόλι την υπογράφει με το κλειδί του χρήστη και την
υποβάλλει. Το ιδιωτικό κλειδί δεν γίνεται ποτέ προσβάσιμο στην εφαρμογή.

Συνέπεια της αρχιτεκτονικής είναι ότι η εφαρμογή δεν αποτελεί σημείο εμπιστοσύνης. Ένας
χρήστης που δεν την εμπιστεύεται μπορεί να καλέσει το συμβόλαιο απευθείας από τη γραμμή
εντολών ή να χρησιμοποιήσει διαφορετική διεπαφή. Όλες οι εγγυήσεις των ΜΛΑ1 έως ΜΛΑ7
παρέχονται από το συμβόλαιο.

\section{Σχεδίαση διεπαφής}
\label{sec:ui-design}

Η διεπαφή σχεδιάσθηκε εκ νέου σε δεύτερη φάση της εργασίας, αφού η πρώτη υλοποίηση, αν και
λειτουργικά πλήρης, είχε τη μορφή γενικού πίνακα διαχείρισης και δεν ανεδείκνυε τη φύση
της συναλλαγής. Οι αρχές του επανασχεδιασμού ήταν οι εξής.

\begin{description}[leftmargin=1.2em,itemsep=2pt]
  \item[«Δίνεις \swap{} παίρνεις».] Κάθε προσφορά παρουσιάζεται, σε όλες τις οθόνες, ως δύο
    πλευρές μιας ανταλλαγής από την οπτική του χρήστη που την κοιτάζει, και ποτέ ως φόρμα ή
    ως γραμμή πίνακα. Η ίδια προσφορά διαβάζεται αντίστροφα από τον προσφέροντα και από
    τον κάτοχο (ΛΑ10).
  \item[Χρώμα με σημασία.] Τρία χρώματα έχουν σταθερό ρόλο: ένα για τις ενέργειες του
    χρήστη, ένα για ό,τι αφορά προσφορές και ανταλλαγές, και ένα για ό,τι παραπέμπει στην
    αλυσίδα (διευθύνσεις, σύνδεσμοι προς τον εξερευνητή).
  \item[Σπανιότητα ως \en{badge}.] Οι τρεις βαθμίδες αποδίδονται ως χαλκός, άργυρος και
    χρυσός, σε τρεις μορφές (ετικέτα, κουκκίδα, πλαίσιο), ώστε να διαβάζονται σε κάθε
    μέγεθος, από την κάρτα έως τη μικρογραφία.
  \item[Σύνθεση προσφοράς σε βήματα.] Η υποβολή γίνεται σε τρία βήματα πλήρους οθόνης
    (\en{SUI}, \en{NFTs}, επισκόπηση), με τη σύνοψη της ανταλλαγής μόνιμα ορατή.
  \item[Ορατή εκκρεμότητα.] Κάθε συναλλαγή διαρκεί μερικά δευτερόλεπτα. Το κουμπί που την
    εκκίνησε δείχνει την πρόοδο, και οι ενέργειες που θα συγκρούονταν μαζί της είτε
    απενεργοποιούνται έως ότου ολοκληρωθεί είτε, στη δημιουργία \en{NFT}, μπαίνουν σε ουρά.
    Κάθε ολοκληρωμένη συναλλαγή συνοδεύεται από σύνδεσμο προς τον εξερευνητή της αλυσίδας.
\end{description}

Το Σχήμα~\ref{fig:wireflow} δείχνει τις οθόνες και τις μεταβάσεις μεταξύ τους. Οι
μακέτες υψηλής πιστότητας σχεδιάστηκαν ως διαδραστικό πρωτότυπο \en{HTML} για
επιτραπέζια και κινητή οθόνη και υλοποιήθηκαν όπως περιγράφεται στο
Κεφάλαιο~\ref{ch:dapp}. Οι τελικές οθόνες, σε πλάτη 1280 και 390 εικονοστοιχείων,
παρουσιάζονται στο Κεφάλαιο~\ref{ch:epideixi}.

\begin{figure}[ht]
\centering
\begin{tikzpicture}[
  font=\footnotesize, >=Stealth,
  scr/.style={draw, rounded corners=3pt, minimum width=2.9cm, minimum height=0.95cm, align=center, fill=gray!8},
  dlg/.style={draw, rounded corners=3pt, minimum width=2.9cm, minimum height=0.95cm, align=center, fill=white, densely dashed}]
  \node[scr] (home) at (0,0) {\en{Marketplace}\\{\scriptsize\code{/}}};
  \node[scr] (coll) at (4.9,0) {Συλλογή\\{\scriptsize\code{/marketplace/:c?rarity=}}};
  \node[dlg] (bid) at (10.2,0) {Υποβολή προσφοράς\\{\scriptsize 3 βήματα, το πρώτο προαιρετικό}};
  \node[scr] (mint) at (0,-2.0) {\en{Mint}\\{\scriptsize\code{/mint}}};
  \node[scr] (mynft) at (4.9,-2.0) {\en{My NFTs}\\{\scriptsize\code{/my-nfts}}};
  \node[dlg] (inc) at (10.2,-2.0) {Εισερχόμενες προσφορές /\\δεσμευμένο \en{NFT}};
  \node[scr] (mybids) at (4.9,-4.0) {\en{My bids}\\{\scriptsize\code{/my-bids}}};
  \node[dlg] (tabs) at (10.2,-4.0) {\en{Received} / \en{Placed}\\{\scriptsize αποδοχή, απόρριψη, ακύρωση}};
  \draw[->] (home) -- (coll);
  \draw[->] (coll) -- node[above, font=\scriptsize] {\en{Place a bid}} (bid);
  \draw[->] (home) -- (mint);
  \draw[->] (coll.south west) -- (mint.north east);
  \draw[->] (mynft) -- node[above, font=\scriptsize] {επιλογή \en{NFT}} (inc);
  \draw[->] (mybids) -- (tabs);
  \node[font=\scriptsize, align=center] at (0,-3.8) {Οι τέσσερις σελίδες είναι\\προσβάσιμες από τη\\γραμμή πλοήγησης};
\end{tikzpicture}
\caption{Οθόνες της εφαρμογής και μεταβάσεις. Με συνεχές περίγραμμα οι σελίδες, με
διακεκομμένο τα παράθυρα διαλόγου.}
\label{fig:wireflow}
\end{figure}
